\paragraph{Checking canonical prime lists.}
Represent a pattern at a source by its set of postings. At an observer,
let $L$ be the ground set of permitted postings. Call $T\subseteq L$
consistent if it contains at most one posting at each input. A guard
$G\subseteq L$ is enabled at $T$ when $G\cap T=\varnothing$. For a
listed family $\Gamma$, write $\operatorname{Tr}(\Gamma)$ for its
inclusion-minimal hitting sets, allowing every subset of $L$ in the
definition. A loss is a consistent set at which no guard is enabled.

The following rule is a restricted form of consensus for
multiple-valued cubes \cite[Sec.~II]{RudellSangiovanniVincentelli1987}.
If $a\in G$ and $b\in H$ post different owners at the same input,
resolve them to $(G\setminus\{a\})\cup(H\setminus\{b\})$. Any
consistent pattern enabling the result must enable $G$ or $H$.

\begin{theorem}[Consensus and minimal losses]\label{core:thm:guard-consensus}
The following are equivalent for a list $\Gamma$ of exclusion guards.
\begin{enumerate}
\item $\Gamma$ consists of all prime sufficient guards for its predicate.
\item $\Gamma$ is an inclusion antichain, and every pairwise resolvent
contains some listed guard.
\item $\Gamma$ is an inclusion antichain, and every member of
$\operatorname{Tr}(\Gamma)$ is consistent.
\end{enumerate}
When these conditions hold, $\operatorname{Tr}(\Gamma)$ is exactly the
set of minimal losses. In particular, completeness of the prime list
can be checked in polynomial time, also when the number of owners varies.
\end{theorem}
\begin{proof}
Every resolvent is sufficient by the rule, so (1) implies (2).
For (2) implies (3), suppose a minimal hitting set $T$ contains competing
postings $a,b$. Each has a private guard, say $G\cap T=\{a\}$ and
$H\cap T=\{b\}$, since deleting either destroys the hitting property.
Their resolvent is disjoint from $T$ and therefore cannot contain a
listed guard.

For (3) implies (1), a guard $B$ containing no listed guard has
$L\setminus B$ meeting every guard. Shrink $L\setminus B$ to a minimal
hitting set. That set is consistent, so it is a loss enabling $B$.
Thus the sufficient guards are exactly the supersets of listed
guards, and the antichain is the prime list. Consistency is inherited
by subsets, so minimal consistent hitting sets are also minimal among
all hitting sets. For the time bound, enumerate pairs of listed guards
and pairs of competing postings, testing inclusion against the same list.
\end{proof}

A restriction that forbids some postings is computed by dropping
those literals from each guard and keeping the inclusion-minimal results.
Indeed a sufficient guard on the restricted domain becomes sufficient
on the original domain after adjoining all the forbidden postings.
It must then contain an original prime, whose restriction it contains.
Consequently two canonical observer lists agree on their overlap
exactly when forbidding both observer colors gives the same reduced list.

\begin{corollary}[One prime list per source]\label{core:cor:common-prime-list}
Fix a source $v$ and canonical lists $\Gamma_{k,v}$ agreeing on the
observer overlaps. Use the ground set of postings at all owners, with
observer lists embedded by inclusion. The primes sufficient on the
proper colorings are exactly the inclusion-minimal sets in
\[
 \widehat\Gamma_v=\left\{\bigcup_{k\in I}G_k:
                             G_k\in\Gamma_{k,v}\text{ for all }k\right\}.
\]
The ordinary minimal hitting sets of $\widehat\Gamma_v$ are exactly the
minimal proper losses. Thus the list also describes the greatest
antitone extension to arbitrary sets of postings, even if they use
all colors or post several colors at one input. At fixed $K$ the common
list and the overlap check take polynomial time in the facet lists.
\end{corollary}
\begin{proof}
A guard is sufficient on every proper coloring precisely when its
restriction to each observer facet contains a listed guard there.
Equivalently it contains a union in the display. On a proper coloring,
success at one facet implies success at the other facets after
deleting their hidden color, by antitonicity and overlap agreement.
Hence a proper coloring loses for the common predicate precisely when
it is a loss in a facet containing it.

For an arbitrary posting set $T$, hitting every union in the display
is equivalent to hitting every guard in at least one facet list.
Indeed, avoiding one guard in each facet gives a union disjoint from
$T$, and the converse is immediate. If $T$ hits an entire facet list,
it contains an ordinary minimal transversal of that facet, which is
consistent by \cref{core:thm:guard-consensus} and omits that observer's
color. This proves the minimal-loss claim. Any antitone extension is
false above every proper loss. The common list is false exactly there,
so it gives the greatest extension. Enumerating the unions costs at
most a product of $K$ list sizes, and removing supersets and comparing
the restricted lists are polynomial operations.
\end{proof}

The greatest extension can also be used for upper bounds on occurrence
filters. An antitone predicate of independent posting sets lies below
the common list everywhere if and only if it lies below it on proper
colorings. Indeed, any false bit of the common list contains a proper
loss, and a predicate omitting that loss also omits its supersets.
